\chapter{Introduzione}

\begin{comment}
Alcune linee guida per scrivere la tesi: 1) Ogni immagine/tabella deve avere le seguenti caratteristiche: a) Una caption; b) Deve essere esplicitamente commentata all'interno del testo (ricordatevi che se vi riferite alla figura/tabella inserendo anche il corrispettivo numero e.g., Figura X mostra che... allora la prima lettera deve essere scritta in maiuscolo; E ricordatevi che latex mette a disposizione un comando per citare correttamente (ref) senza dover mettere a mano il numeretto corrispondente) c) Le immagini devono essere rifatte tramite PPT o similari e salvate in pdf prima di inserirle nel file latex, altrimenti rischiate di inserire immagini sgranate. 2) Ricordatevi sempre di citare le fonti e di rifrasare, altrimenti rischiate di inccorrere nel reato di plagio (che può comportare anche la revoca del titolo accademico) e che citare non sminuisce il vostro lavoro, anzi, lo rafforza. 3)L'Abstract deve necessariamente contenere le seuguenti informazioni a) Un breve riasssunto della problematica; b)Una breve descrizione di ciò che è stato fatto; c) I risultati ottenuti.
 
\end{comment}

\section{Contesto e motivazione}

I sistemi che integrano componenti di machine learning (\textit{ML-enabled systems}) sono ormai parte del software di uso quotidiano e con la loro diffusione è cresciuta l'attenzione per la qualità del codice che li realizza. \cite{hidden_tech_debt} hanno mostrato che questi sistemi accumulano debito tecnico con particolare facilità, perché al codice tradizionale si aggiungono dipendenze dai dati, pipeline di addestramento e configurazioni difficili da tenere sotto controllo. Una delle forme più studiate di debito tecnico è il code smell: un sintomo, nel codice, di una scelta di progettazione o di implementazione discutibile, che non impedisce al programma di funzionare ma ne rende più costosa la manutenzione.
Per il codice di machine learning, \cite{code_smells_definitions} hanno proposto un catalogo di smell specifici (ML-specific code smell, nel seguito anche ML-CS o semplicemente smell), ricavato dalla letteratura scientifica e dalla documentazione tecnica: usi impropri delle API di Pandas, NumPy, TensorFlow e PyTorch che possono compromettere prestazioni, correttezza e riproducibilità di una pipeline. \cite{ml_specific_gib} hanno poi studiato il ciclo di vita di questi smell su larga scala. Il loro lavoro introduce CodeSmile, uno strumento di analisi statica che rileva dodici ML-CS, e lo applica a 337 progetti e a oltre 400.000 commit del dataset NICHE \cite{niche}, ricostruendo per ogni smell il commit che lo introduce e, quando esiste, il commit che lo rimuove.
Per capire in quali attività gli smell vengono introdotti e rimossi, lo studio classifica i commit in quattro categorie — \textbf{nuova funzionalità}, \textbf{correzione di bug}, \textbf{miglioramento e refactoring} — applicando al messaggio di commit la ricerca di parole chiave già adottata da \cite{smells_bad} per gli smell tradizionali. Tra i risultati, due riguardano direttamente il refactoring: una parte consistente delle introduzioni avviene in commit etichettati come refactoring, miglioramenti e refactoring sono le attività più frequenti nei commit di rimozione.
Questa classificazione ha però un limite, da cui nasce la presente tesi. L'etichetta «refactoring» dipende soltanto da ciò che lo sviluppatore ha scritto nel messaggio di commit: non verifica che nel codice sia avvenuta davvero un'operazione di refactoring, né che questa abbia toccato il punto in cui si trova lo smell. Un commit descritto come \textit{refactoring} può non contenere alcuna ristrutturazione riconoscibile; al contrario, un commit che rinomina o estrae un metodo può avere un messaggio che non lo dice. Gli stessi autori di CodeSmile osservano che alcune rimozioni sono probabilmente incidentali, cioè conseguenza di modifiche strutturali come la ridenominazione di un file, e che lo spostamento di uno smell viene registrato come una rimozione seguita da una nuova introduzione.
La domanda che resta aperta è quindi la seguente: \textit{quando uno smell di machine learning compare o scompare, nel codice di quel commit è avvenuto davvero un refactoring e riguarda proprio la funzione che contiene lo smell?} Per rispondere servono uno strumento che rilevi i refactoring dalle modifiche al codice e un criterio che colleghi ciascun refactoring alla posizione dello smell. In questa tesi lo strumento è PyRef \cite{pyref}, che riconosce i refactoring a livello di metodo nei progetti Python ed il criterio è la coincidenza di \textbf{file e funzione} tra il refactoring e lo smell.

/section{Obiettivo e domande di ricerca}


/section{Requisiti dello studio}

/section{Contributi}

/section{Struttura della tesi}